%!TEX root = main.tex

\section{Evaluated Dissemination Algorithms}
\label{sec:algorithms}


In general, dissemination methods share a set of steps that include ``\textit{variation points}''.
The methods are then classified based on the concrete way in which these variation points are implemented.
For instance, the foremost variation point is the usage or not of an \textit{underlying topology} to drive the dissemination process.

On the other hand, methods that do not build an underlying topology can be classified using additional features.
For instance, how much contextual information they use.
In this case, we find \textit{context-aware} and \textit{context-oblivious} methods, but there are also methods where nodes only use a specific piece of contextual information (e.g., neighbor knowledge).
The algorithms of this group that we consider in this paper are context-aware and \textit{statistical}.
They share the following common structure, where $T$ is a value related to the context and $T_C$ is an associated threshold:

\begin{enumerate}
	\item \textbf{Initialize} $T$, \textbf{calculate threshold} value $T_C$.
	\item When a broadcast message $m_i$ is received, \textbf{update the value} of $T$.
		  
		  If $T < T_C$, do not retransmit and stop procedure. Otherwise, set random timer.
		  
	\item If $m_i$ is received again before the timer expires, repeat the previous step.
	\item When the timer expires, retransmit $m_i$ if $T > T_C$.
\end{enumerate}

\subsection{Without underlying topology}

The \textbf{distance-to-mean} (dist2mean) approach~\cite{Slavik2011} is based on a simple observation: if a node $v$ receives a message from all directions, there is little benefit in retransmitting it because likely all neighbors already received the message.
In other words, the neighborhood of $v$ is already covered.
To estimate the coverage of $v$, the method computes a metric for each broadcast session -- the distance to the spatial mean.
This is, every time the node $v$ receives a message, it records the position $(x_i, y_i)$ of the transmitting node.
When the timer expires and it is time to decide whether retransmit, this only happens if $T > T_C$, where:

\begin{equation}
T = \dfrac{\sqrt{(v_x-\bar{x})^2 + (v_y-\bar{y})^2}}{r}
\end{equation}

\begin{equation}\label{eq:dist2mean_threshold}
T_C = 0.95 -0.7e^{-0.11N}
\end{equation}

\begin{equation}
M = (\bar{x},\bar{y}) = (\dfrac{1}{n}\sum_{i=1}^{n}x_i, \dfrac{1}{n}\sum_{i=1}^{n}y_i)
\end{equation}

and $N$ is the number of neighbors of $v$ and $(v_x, v_y)$ is its location.
The value of function~\ref{eq:dist2mean_threshold}, which have been taken from~\cite{Slavik2011}, decreases as $N$ decreases in order to ensure connectivity.
